\documentclass[11pt]{article}
\usepackage[margin=1in]{geometry}
\usepackage{amsmath,amssymb}
\usepackage{booktabs}
\usepackage{tabularx}
\usepackage{array}
\usepackage{xcolor}
\usepackage[colorlinks=true,linkcolor=blue]{hyperref}

% badge to flag effects that are new in 3D (absent from the 2D model)
\newcommand{\newin}{{\small\textbf{\color{blue}[3D]}}}
\newcommand{\drop}{{\small\textbf{\color{red!70!black}[omitted]}}}

\newcolumntype{L}{>{\raggedright\arraybackslash}X}

\title{\vspace{-1.5em}3D Seed Dynamics: Physics of the Strip-Theory Model}
\author{}
\date{}

\begin{document}
\maketitle
\vspace{-3em}

This note summarizes the physics carried over from the 2D Andersen--Pesavento--Wang
model (\texttt{minimal\_imp}) into the 3D strip-theory model, what each piece does
physically, and where it lands in the code. The \newin{} badge marks effects that
are genuinely new in 3D; \drop{} marks 2D terms we deliberately drop (all negligible
for a seed in air and exactly zero for the symmetric test seed).

\section*{Setup and conventions}
The plate is sliced into spanwise \emph{strips}; each strip is a local 2D
cross-section using the exact 2D coefficient laws. Body axes: $x$ = chordwise,
$y$ = plate-normal, $z$ = spanwise. Everything is referenced to the center of
mass (CoM). State: $[\,\mathbf{r}\ (\text{CoM, inertial});\ \mathbf{q}\ (\text{body}\to\text{world});\ \mathbf{v}\ (\text{inertial});\ \boldsymbol{\omega}\ (\text{body})\,]$.
Per strip we use the in-plane relative velocity $\mathbf{v}_{ip}=(v_x,v_y)$ with
speed $|v_{ip}|$, the local angle of attack $\alpha=\operatorname{atan2}(v_y,v_x)$,
and the spanwise spin component $\omega_z$. Coefficients $C_T,C_D,C_R,C_{D2}$ come
from \texttt{computeAeroCoeffs}.

\section*{Per-strip forces}
Each strip of chord $c$ and width $dz$ contributes the following. Lift is
$\perp$ to the relative velocity, drag is $\parallel$ to $-\mathbf{v}$. Forces are
summed over strips; gravity is added once at the end.

\renewcommand{\arraystretch}{1.4}
\begin{tabularx}{\textwidth}{@{} l >{\raggedright\arraybackslash}p{0.34\textwidth} L @{}}
\toprule
\textbf{Force} & \textbf{Rough magnitude \& direction} & \textbf{Acts at} \\
\midrule
Translational lift &
$\tfrac12\,\rho\,c\,dz\,C_T\,|v_{ip}|^2$, \ $\perp \mathbf{v}$ &
center of pressure, $x_{cp}=l_{cp}\cdot c$ (migrates with $\alpha$) \\

Translational drag &
$\tfrac12\,\rho\,c\,dz\,C_D\,|v_{ip}|^2$, \ $\parallel -\mathbf{v}$ &
center of pressure (same as lift) \\

Rotational lift \newin{} &
$\tfrac12\,\rho\,c^2\,dz\,C_R\,\omega_z\,|v_{ip}|$, \ $\perp \mathbf{v}$ &
mid-chord ($l_{cr}=0$) \\
\bottomrule
\end{tabularx}

\medskip
\noindent\textbf{Translational lift} — circulatory force from the flow turning
around the plate; the workhorse that holds the seed up.\\
\textbf{Translational drag} — pressure/separation resistance opposing motion.\\
\textbf{Rotational lift} — Magnus-like lift from the plate's own spin shedding
circulation. \newin{} In 3D only the spanwise spin $\omega_z$ drives it (the
in-plane rotation each strip feels), and it acts at mid-chord, not the CoP.

\section*{Per-strip torques (about the CoM)}
In 3D a force at position $\mathbf{r}$ gives torque $(\mathbf{r}-\mathbf{c})\times\mathbf{F}$
automatically---the cross product reproduces the 2D moment arm
$l_\tau=l_{cp}-l_{cm}$ with no special handling.

\begin{tabularx}{\textwidth}{@{} l >{\raggedright\arraybackslash}p{0.36\textwidth} L @{}}
\toprule
\textbf{Torque} & \textbf{Rough equation} & \textbf{Arm / axis} \\
\midrule
Lift+drag moment ($T_t$) &
$(\mathbf{r}_{cp}-\mathbf{c})\times \mathbf{F}_{\text{lift+drag}}$ &
arm = CoP $-$ CoM \\

Rotational-lift moment ($T_{cr}$) &
$(\mathbf{r}_{lcr}-\mathbf{c})\times \mathbf{F}_{\text{rot lift}}$ &
arm = mid-chord $-$ CoM \\

Spin damping ($T_r$) &
$-\tfrac{1}{128}\,\rho\,c^4\,dz\,C_{D2}\,\omega_z|\omega_z|\cdot(\text{end factor})$ &
about spanwise ($z$) axis \\

Buoyancy couple ($T_b$) \drop{} &
$-m_p\,g\,l_{cm}\cos\theta$ &
weight at CoM vs.\ buoyancy at mid-chord \\
\bottomrule
\end{tabularx}

\medskip
\noindent\textbf{$T_t$} — the pitching torque from lift+drag acting ahead of/behind
the CoM; the dominant driver of flutter vs.\ tumble.\\
\textbf{$T_{cr}$} — moment of the rotational lift; kept separate because it acts at
a different chordwise point ($l_{cr}=0$) than $T_t$.\\
\textbf{$T_r$} — nonlinear spin damping: each chordwise element sees drag
$\propto(\omega r)^2$, and integrating $r\times$drag along the chord opposes spin.
This is \emph{separate physics}, not captured by the strip lift/drag (which are
evaluated at mid-chord where $\omega_z$ contributes no velocity). \newin{} Taken
about the spanwise axis only; damping about the other axes is captured implicitly
through the $\boldsymbol\omega\times\mathbf r$ term in each strip's velocity.\\
\textbf{$T_b$} \drop{} — pendulum-like restoring couple from weight and buoyancy
acting at different points. Negligible in air ($m_p=\rho_{\text{air}}V\ll m$).

\section*{Added mass}
Fluid the plate must accelerate with it; scales with fluid density, so small in air
but retained. Built per strip and summed (\texttt{getAddedMass}).

\begin{tabularx}{\textwidth}{@{} l L @{}}
\toprule
\textbf{Block} & \textbf{Content} \\
\midrule
Translational $A_{\text{trans}}$ &
$\operatorname{diag}\!\big(0,\ \textstyle\sum \pi\rho(c/2)^2 dz,\ 0\big)$: only broadside (normal, $y$)
entrains fluid; edge-on ($x$) and in-plane span ($z$) $\approx 0$. \\
Rotational $A_{\text{rot}}$ &
$I_{zz}=\sum(i_a+m_2 x_s^2)$ (the 2D $I_a$); \ \ $I_{xx}=\sum m_2 z_s^2$ \newin{}; \ \
$I_{yy}\approx 0$ \newin{}; \ \ $I_{xz}=-\sum m_2 x_s z_s$ (zero if symmetric). \\
\bottomrule
\end{tabularx}

\noindent Here $m_2=\pi\rho(c/2)^2 dz$, $i_a=\pi\rho c^4 dz/128$, and $x_s,z_s$ are
the strip's chordwise/spanwise offsets from the CoM.

\section*{What changes between 2D and 3D}

\paragraph{Reference frame / Coriolis terms.} The 2D code works in the body frame,
so it carries rotating-frame terms $(m{+}m_2)\omega v_{yp}$, $-(m{+}m_1)\omega v_{xp}$.
We sum translational forces in the \emph{inertial} frame, where these vanish
identically---provided $A_{\text{trans}}$ is rotated into the inertial frame,
$R\,A_{\text{trans}}R^{\top}$, before dividing.

\paragraph{Added-mass CoM-offset coupling. \drop{}} The 2D terms
$-m_2\omega^2 l_{cm}$ and $+m_2\dot\omega\, l_{cm}$ are the translation$\leftrightarrow$rotation
added-mass coupling (fluid centered at mid-chord, CoM offset). We neglect the
added-mass cross-block, dropping these. Justified: zero for the symmetric seed
($l_{cm}=0$) and $\propto\rho_{\text{air}}$ otherwise.

\paragraph{Gyroscopic torque. \newin{}} The rotational equation gains
$\boldsymbol\omega\times(I\boldsymbol\omega)$, which is identically zero for scalar
2D rotation but real in 3D (spinning about one axis induces torque about others).

\paragraph{Time-varying inertia $\dot I_G$. \newin{}} If the CoM drifts (e.g.\ a
shifting nut), $I_G$ changes in time, adding $\dot I_G\,\boldsymbol\omega$ to the
rotational equation. 2D inertia was constant.

\paragraph{Velocity sampling point.} 2D evaluates the reference velocity at the
CoM ($v_{yp}-\omega l_{cm}$); per strip we evaluate at the strip mid-chord and let
$\boldsymbol\omega\times\mathbf r$ supply the rotational part. This is why $T_r$ must
stay explicit---mid-chord sampling misses the chordwise spin-drag distribution.

\paragraph{Attitude. \newin{}} A full quaternion replaces the single angle
$\theta$; orientation advances by quaternion kinematics
$\dot{\mathbf q}=\tfrac12\,\mathbf q\otimes[0,\boldsymbol\omega]$.

\paragraph{Effective weight.} 2D uses $(m-m_p)g$ (weight minus buoyancy); we use
$m\,g$, consistent with dropping buoyancy.

\section*{Where it fits in the implementation}
\begin{tabularx}{\textwidth}{@{} l L @{}}
\toprule
\textbf{Function} & \textbf{Role} \\
\midrule
\texttt{getMassProperties} & CoM $\mathbf c$, $I_G$, $\dot I_G$, and added mass at time $t$ \\
\texttt{getAddedMass} & $A_{\text{trans}}$, $A_{\text{rot}}$ (tables above) \\
\texttt{computeSeedLocalVel} & per-strip body-frame velocity via $\mathbf v+\boldsymbol\omega\times\mathbf r$ \\
\texttt{computeAeroCoeffs} & $C_T,C_D,l_{cp},C_R,C_{D2}$ vs.\ $\alpha$ \\
\texttt{stripForce} \emph{(to write)} & the three forces in the force table \\
\texttt{stripSpinDamping} \emph{(to write)} & the $T_r$ term \\
\texttt{quatKinematics} \emph{(to write)} & $\dot{\mathbf q}$ update \\
translation / rotation EOM & sum forces (+gravity) $\to \mathbf a$; sum torques $+\,\boldsymbol\omega\times(I\boldsymbol\omega)+\dot I_G\boldsymbol\omega \to \boldsymbol\alpha$ \\
\bottomrule
\end{tabularx}

\end{document}
